\chapter{Article I : Titre} \label{chap:article1}

{\large \textbf{Description de l’article}}

\begin{itemize}
  \item \textbf{Titre :} Analysis Spatialis Migrationis Larorum in Peruvia inter Annos 1942 et 1950 : El Niño, Phenomenon Exogena vel Non?
  
  \item \textbf{Nom de la revue à comité de lecture :} Geographical Analysis
  
  \item \textbf{Auteurs.trices :}
    \begin{itemize}
    \item Prénom \textsc{Nom}, Département de géomatique appliquée, Université de Sherbrooke.
    \item Prénom \textsc{Nom}, Département à compléter, Université à compléter.
    \item Prénom \textsc{Nom}, Département à compléter, Université à compléter.
  \end{itemize}
  \item \textbf{Référence complète de l'article :} Inclure le DOI si l'article est déjà publié.

  \item \textbf{État de cheminement :}
        \begin{itemize}
            \item $\checkmark$ Publié le jour mois année
            \item $\square$ Sous presse
            \item $\square$ Accepté moyennant modifications
            \item $\square$ Soumis le jour mois année
   \end{itemize}


  
  \item \textbf{Résumé :} \begin{spacing}{.8} 
  \small Estimatio Densitatis Nucleorum (KDE) est methodus late adhibita in geographia ad studium concentrationis punctorum distributionis. Retia geographica sunt spatia unius et dimidiae dimensionis, quae proprietates specificas exhibent, cum res in his retibus eveniunt (accidentia in viis, ruptiones canalium, species iuxta flumina, etc.). Hisce in annis decem proximis, haec natura reticulata postulavit extensionem KDE spatii. Variae versione KDE Reticulatae (NKDE) propositae sunt, quarum unaquaeque suos habet usus ac limitationes, atque nunc regulariter adhibentur. Nihilominus, parva cura adhibita est extensioni temporali NKDE, scilicet TNKDE. In praxi, cum eventa studenda temporibus definitis accidant et retibus constringantur, geographi methodologias utuntur quae vel dimensionem reticularem vel temporalem negligere solent. Hic propositum est methodum TNKDE innixam recentioribus progressibus NKDE necnon prodotto nucleorum. Simul adaptamus methodos classicas KDE, ut correctio Diggle, latitudo adaptiva Abramson et selectio latitudinis per verisimilitudinem maximam "leave-one-out". Methodus quoque illustratur per exemplum accidentium viariorum in Monte Regio (Montreal), quae pedestres implicaverunt inter annos 2016 et 2019.
  \end{spacing}
\end{itemize}

{\large \textbf{Contributions de chaque coautrice et coauteur}}
\begin{itemize}
  \item \textbf{Auteur.trice I} 
  \begin{spacing}{1}  \small 
  Conceptualisation; revue de littérature; méthodologie; collecte de données; analyses de données; interprétation des résultats; rédaction – première ébauche; rédaction – révision et édition; visualisation (figures, graphiques, tableaux, etc.); supervision; acquisition des fonds; gestion du projet. 
  \end{spacing}

  \item \textbf{Auteur.trice II} 
      \begin{spacing}{1}  \small 
      Conceptualisation... 
      \end{spacing}
  
  \item \textbf{Auteur.trice III} 
    \begin{spacing}{1}  \small 
    Conceptualisation... 
    \end{spacing}
\end{itemize}

\newpage
\section{Section 1} \label{sec:titre_articlesection1}

\textbf{Attention,} si votre article est en anglais, changez la langue avec \texttt{\textbackslash selectlanguage\{english\}}. Puis, pour Pour remettre la langue est français, utiliser \texttt{\textbackslash selectlanguage\{french\}}.

\lipsum[1-4] % génère des paragraphes (à supprimer)


\section{Section 2} \label{sec:titre_articlesection2}

Paragraphe.

\lipsum[1-4] % génère des paragraphes (à supprimer)


